\chapter{曲线运动}
\label{chap:curvilinear-motion}

\begin{learninggoals}
  \begin{itemize}[leftmargin=2em]
    \item 判断物体做曲线运动的条件；
    \item 理解速度方向与轨迹切线的关系；
    \item 为平抛运动和圆周运动的分解方法做准备。
  \end{itemize}
\end{learninggoals}

\sourcetodo{本章根据“物理/物理1.pdf”开始整理。该文件共 20 页，目前仅抽样核对了封面和第 2 页。}

\section{曲线运动的基本特征}

物体的运动轨迹为曲线时，称物体做曲线运动。质点在某一点的瞬时速度方向沿轨迹在该点的切线方向，因此曲线运动的速度方向不断变化。

\begin{definitionbox}[曲线运动的条件]
  当物体所受合外力的方向与速度方向不在同一直线上时，速度方向将发生变化，物体可能做曲线运动。
\end{definitionbox}

\begin{examplebox}[运动的合成]
  研究复杂运动时，可以根据相互垂直的坐标轴把位移、速度和加速度分解，再分别研究各分运动。
\end{examplebox}

\begin{exercisebox}
  一个物体速度向东，而所受恒力向北。只讨论速度方向的变化，说明它为什么不会继续做直线运动。
\end{exercisebox}
